\documentclass[12pt,a4paper]{article}
\usepackage[UTF8]{ctex}
\usepackage{amsmath}
\usepackage{graphicx}
\usepackage{booktabs}
\usepackage{geometry}
\usepackage{siunitx}
\usepackage{pgfplots}
\usepackage{ulem}
\usepackage{caption}
\pgfplotsset{compat=1.17}
\geometry{left=2.5cm,right=2.5cm,top=2.5cm,bottom=2.5cm}
\sisetup{per-mode=symbol}

\title{\Huge{\textbf{\uuline{实\ 验\ 报\ 告}}}}
\author{\uline{少年班学院} \quad 姓名: \uline{李大恕} \quad 学号: \uline{PB24000145}}
\date{\today}

\begin{document}
\maketitle

\section{实验题目}
电流电离室的饱和特性和轴向均匀性测量

\section{实验目的}
\begin{enumerate}
    \item 了解电流电离室的基本原理和饱和特性。
    \item 了解弱电流测量与电压-频率（V/F）变换的基本方法。
    \item 学习用电流电离室进行电离辐射探测与轴向均匀性测量。
\end{enumerate}

\section{实验原理}
\subsection{电流电离室与饱和区}
带电粒子或中性粒子在电离室灵敏体积内直接或间接引起气体电离，在外加电场作用下形成电离电流。极间电压 $V$ 较低时，正负离子复合较多，收集不完全；随 $V$ 升高，复合减小，电流增大；当电压升至一定值 $V_0$ 后，电离电荷几乎被全部收集，进入饱和区，对应饱和电流 $I_0$。

若单位时间进入电离室的粒子数为 $n$，单个粒子在气体中损失平均能量为 $E$，产生一对离子所需能量为 $\epsilon$，则饱和电流为
\begin{equation}
    I_0 = n\frac{E\,\mathrm{e}}{\epsilon},
\end{equation}
其中 $\mathrm{e}$ 为元电荷。

\subsection{弱电流放大与读出}
实验讲义中的读出链路为“弱电流放大器 + V/F 变换 + 定标器”。忽略失调和噪声时，放大器输出电压 $V_i$ 与输入电流 $I_i$ 成正比：
\begin{equation}
    V_i = -R_L I_i.
\end{equation}
V/F 变换输出频率满足
\begin{equation}
    f_i = G V_i,
\end{equation}
故定标器计数率可作为电离电流的线性表征量。

\subsection{轴向均匀性}
在选定饱和工作电压下，将准直后的 $^{137}\mathrm{Cs}$ $\gamma$ 源沿电离室轴向平移，记录不同轴向位置的读数，得到轴向响应曲线。若轴向均匀性好，则中段响应应接近常数。

\begin{figure}[htbp]
    \centering
    \includegraphics[width=0.78\textwidth]{assets/page4_img1.png}
    \caption{讲义给出的弱电流放大器与输出电路原理图}
\end{figure}

\section{实验步骤（按讲义与本次数据对应）}
\begin{enumerate}
    \item 连接电离室、弱电流放大器、V/F 变换器和定标器，预热并校零。
    \item 在不同辐照强度条件下，扫描极间电压，测量 V-I 饱和曲线。
    \item 在辐照强度不变时，改变电离室充气压力，测量不同压力下的 V-I 曲线。
    \item 在选定工作电压（本次数据为 \SI{300}{\volt}）下，沿轴向移动准直源，测量轴向均匀性。
\end{enumerate}

\section{原始数据与预处理}
\subsection{重复测量与均值定义}
本次 xlsx 中每个测点均有两组测量值。后续处理使用算术平均：
\begin{equation}
    \bar{N}(x)=\frac{N_1(x)+N_2(x)}{2},
\end{equation}
其中 $x$ 表示电压或轴向位置。

\subsection{数据核对项（待确认）}
在“\SI{2.5}{MPa}”条件下，\SI{67.5}{\volt} 点两次记录分别为 1219 与 12221。第二个值与同曲线相邻点量级（约 $10^3$）明显不一致，判断为记录异常。为避免引入人为猜测，本报告在作图时该点采用第一组原始值 1219，并保留该点为“待核对数据”。

\section{数据处理与结果}
\subsection{不同辐照强度下的 V-I 曲线}
将 0 块板、1 块板、2 块板、3 块板四种工况分别取两组均值后作图，如图~\ref{fig:vi-intensity}。可见辐照强度减弱时，平台电流依次降低。

\begin{figure}[htbp]
\centering
\begin{tikzpicture}
\begin{axis}[
    width=0.88\textwidth,
    height=0.44\textwidth,
    xlabel={极间电压 $V$ (\si{\volt})},
    ylabel={计数 $\bar{N}$ (1 s)},
    xmin=0, xmax=450,
    ymin=0, ymax=1600,
    grid=both,
    legend pos=south east,
    mark size=1.3pt
]
\addplot[blue, thick, mark=*] coordinates {(0,86.5) (7.5,551.5) (15,802.0) (22.5,983.5) (30,1088.5) (37.5,1182.5) (45,1241.5) (52.5,1288.0) (60,1314.5) (67.5,1343.5) (75,1361.5) (82.5,1377.5) (90,1393.5) (97.5,1404.0) (105,1414.0) (112.5,1423.0) (120,1427.5) (127.5,1439.0) (135,1441.5) (142.5,1443.5) (150,1452.0) (165,1457.0) (180,1464.0) (195,1472.5) (210,1472.0) (225,1477.5) (240,1486.0) (255,1485.5) (270,1491.0) (285,1490.5) (300,1494.0) (315,1495.5) (330,1499.5) (345,1502.5) (360,1499.5) (375,1506.0) (390,1507.0) (405,1506.0) (420,1509.0) (435,1452.5) (450,1460.5)};
\addplot[red, thick, mark=square*] coordinates {(0,94.0) (7.5,442.5) (15,639.0) (22.5,775.5) (30,871.5) (37.5,929.5) (45,966.0) (52.5,992.0) (60,1013.5) (67.5,1026.5) (75,1042.0) (82.5,1051.5) (90,1061.5) (97.5,1066.0) (105,1074.0) (112.5,1078.5) (120,1082.0) (127.5,1083.5) (135,1090.0) (142.5,1091.0) (150,1094.0) (165,1097.5) (180,1100.5) (195,1105.5) (210,1108.0) (225,1113.0) (240,1112.5) (255,1116.5) (270,1115.5) (285,1118.5) (300,1121.0) (315,1121.5) (330,1122.5) (345,1126.5) (360,1125.0) (375,1129.0) (390,1128.5) (405,1129.0) (420,1129.5) (435,1135.0) (450,1130.5)};
\addplot[teal!70!black, thick, mark=triangle*] coordinates {(0,47.0) (7.5,405.0) (15,586.5) (22.5,675.5) (30,718.5) (37.5,761.0) (45,785.0) (52.5,801.5) (60,815.0) (67.5,825.0) (75,832.5) (82.5,838.5) (90,843.5) (97.5,847.0) (105,851.0) (112.5,853.0) (120,856.0) (127.5,859.0) (135,862.5) (142.5,862.0) (150,864.5) (165,869.0) (180,871.5) (195,874.5) (210,877.0) (225,879.5) (240,880.5) (255,877.5) (270,885.0) (285,882.0) (300,884.5) (315,885.5) (330,885.5) (345,886.5) (360,890.0) (375,889.5) (390,893.0) (405,892.0) (420,890.0) (435,894.0) (450,892.5)};
\addplot[violet!70!black, thick, mark=diamond*] coordinates {(0,50.0) (7.5,352.5) (15,467.5) (22.5,553.5) (30,594.5) (37.5,617.5) (45,635.5) (52.5,644.5) (60,642.0) (67.5,648.5) (75,652.5) (82.5,658.5) (90,664.0) (97.5,661.5) (105,666.5) (112.5,668.0) (120,668.5) (127.5,671.5) (135,673.0) (142.5,673.5) (150,677.0) (165,678.5) (180,679.0) (195,682.5) (210,683.0) (225,682.5) (240,686.0) (255,685.0) (270,688.0) (285,687.5) (300,693.5) (315,691.0) (330,691.0) (345,691.5) (360,690.5) (375,692.5) (390,693.5) (405,690.5) (420,693.5) (435,695.0) (450,693.5)};
\legend{0块板, 1块板, 2块板, 3块板}
\end{axis}
\end{tikzpicture}
\caption{不同辐照强度下的 V-I 曲线（两组均值）\label{fig:vi-intensity}}
\end{figure}

以 \SI{300}{\volt} 到 \SI{390}{\volt} 区间均值作为平台统计量，得到表~\ref{tab:intensity-summary}。

\begin{table}[htbp]
\centering
\caption{不同辐照强度条件下的平台特征\label{tab:intensity-summary}}
\begin{tabular}{cccc}
\toprule
工况 & $\bar{N}(\SI{300}{\volt})$ & 平台均值 $\bar{N}_{\mathrm{sat}}$ (300--390 V) & 饱和起始电压约值 /V \\
\midrule
0块板 & 1494.0 & 1500.6 & 120 \\
1块板 & 1121.0 & 1124.9 & 105 \\
2块板 & 884.5  & 887.8  & 90 \\
3块板 & 693.5  & 691.9  & 82.5 \\
\bottomrule
\end{tabular}
\end{table}

\subsection{不同充气压力下的 V-I 曲线}
同样按两组均值作图，如图~\ref{fig:vi-pressure}。随压力从 \SI{1}{MPa} 增加到 \SI{2.5}{MPa}，平台计数整体升高，与讲义“压力越高，饱和电流越大、饱和电压越高”的趋势一致。

\begin{figure}[htbp]
\centering
\begin{tikzpicture}
\begin{axis}[
    width=0.88\textwidth,
    height=0.44\textwidth,
    xlabel={极间电压 $V$ (\si{\volt})},
    ylabel={计数 $\bar{N}$ (1 s)},
    xmin=0, xmax=450,
    ymin=0, ymax=1400,
    grid=both,
    legend pos=south east,
    mark size=1.3pt
]
\addplot[black, thick, mark=*] coordinates {(0,112.5) (7.5,379.5) (15,560.0) (22.5,625.0) (30,649.5) (37.5,660.5) (45,665.5) (52.5,673.5) (60,674.5) (67.5,675.5) (75,676.5) (82.5,677.0) (90,677.0) (97.5,677.5) (105,678.5) (112.5,679.0) (120,679.0) (127.5,678.0) (135,678.5) (142.5,677.5) (150,678.5) (165,679.5) (180,679.5) (195,681.5) (210,681.5) (225,681.0) (240,678.0) (255,679.5) (270,679.5) (285,679.5) (300,679.5) (315,681.0) (330,683.0) (345,680.5) (360,681.0) (375,681.0) (390,679.5) (405,681.5) (420,682.0) (435,682.5) (450,681.5)};
\addplot[blue, thick, mark=square*] coordinates {(0,104.0) (7.5,412.5) (15,649.0) (22.5,751.0) (30,810.5) (37.5,827.5) (45,862.5) (52.5,861.0) (60,872.0) (67.5,876.0) (75,880.0) (82.5,884.5) (90,886.0) (97.5,885.5) (105,887.5) (112.5,889.5) (120,892.5) (127.5,892.0) (135,892.5) (142.5,893.5) (150,891.5) (165,892.5) (180,895.0) (195,897.0) (210,896.0) (225,899.0) (240,899.5) (255,899.5) (270,901.0) (285,898.5) (300,898.0) (315,898.0) (330,897.5) (345,899.5) (360,899.5) (375,899.0) (390,901.0) (405,899.0) (420,901.0) (435,903.5) (450,902.5)};
\addplot[red, thick, mark=triangle*] coordinates {(0,101.5) (7.5,442.5) (15,666.0) (22.5,822.0) (30,902.0) (37.5,956.5) (45,988.0) (52.5,1012.5) (60,1025.5) (67.5,1039.0) (75,1045.5) (82.5,1054.5) (90,1059.0) (97.5,1062.5) (105,1068.5) (112.5,1071.0) (120,1070.5) (127.5,1069.5) (135,1073.0) (142.5,1071.5) (150,1075.0) (165,1081.5) (180,1081.5) (195,1083.5) (210,1083.5) (225,1084.5) (240,1085.0) (255,1087.0) (270,1086.5) (285,1087.5) (300,1086.5) (315,1087.0) (330,1087.5) (345,1089.0) (360,1091.5) (375,1087.5) (390,1089.5) (405,1090.0) (420,1091.0) (435,1093.5) (450,1092.0)};
\addplot[teal!70!black, thick, mark=diamond*] coordinates {(0,99.5) (7.5,335.5) (15,646.0) (22.5,838.0) (30,980.0) (37.5,1063.0) (45,1122.5) (52.5,1166.0) (60,1195.5) (67.5,1219.0) (75,1237.0) (82.5,1250.0) (90,1263.0) (97.5,1273.5) (105,1279.0) (112.5,1284.5) (120,1289.5) (127.5,1288.0) (135,1296.0) (142.5,1299.0) (150,1301.0) (165,1303.5) (180,1310.5) (195,1312.5) (210,1313.0) (225,1317.5) (240,1318.0) (255,1324.0) (270,1322.5) (285,1322.0) (300,1326.5) (315,1325.5) (330,1326.5) (345,1331.0) (360,1326.5) (375,1330.0) (390,1331.0) (405,1329.0) (420,1331.5) (435,1334.0) (450,1332.0)};
\legend{1 MPa, 1.5 MPa, 2 MPa, 2.5 MPa}
\end{axis}
\end{tikzpicture}
\caption{不同充气压力下的 V-I 曲线（两组均值）\label{fig:vi-pressure}}
\end{figure}

\begin{table}[htbp]
\centering
\caption{不同压力条件下的平台特征\label{tab:pressure-summary}}
\begin{tabular}{cccc}
\toprule
压力工况 & $\bar{N}(\SI{300}{\volt})$ & 平台均值 $\bar{N}_{\mathrm{sat}}$ (300--390 V) & 饱和起始电压约值 /V \\
\midrule
1 MPa   & 679.5  & 680.8  & 30 \\
1.5 MPa & 898.0  & 898.9  & 45 \\
2 MPa   & 1086.5 & 1088.4 & 67.5 \\
2.5 MPa & 1326.5 & 1328.1 & 67.5$^{\ast}$ \\
\bottomrule
\end{tabular}
\caption*{$^{\ast}$该工况在 \SI{67.5}{\volt} 处存在一条异常记录（12221），已在预处理时标注“待确认”。}
\end{table}

\subsection{轴向均匀性曲线（\SI{300}{\volt}, \SI{10}{s}）}
两组均值后的轴向曲线见图~\ref{fig:axial}。令相对响应
\begin{equation}
    \eta(z)=\frac{\bar{N}(z)}{\bar{N}_{\max}},
\end{equation}
其中 $\bar{N}_{\max}=795.0$（位于 \SI{12}{cm}）。计算得：
\begin{itemize}
    \item $\eta\ge 0.98$ 的均匀区间约为 \SIrange{4}{31}{cm}；
    \item 该区间平均计数为 790.8，非均匀度
    \begin{equation}
        \delta=\frac{\bar{N}_{\max}-\bar{N}_{\min}}{\bar{N}_{\mathrm{avg}}}\times 100\%\approx 1.83\%;
    \end{equation}
    \item 两端响应相对均匀区均值下降：$z=0$ 处约 16.8\%，$z=38$ 处约 23.1\%。
\end{itemize}

\begin{figure}[htbp]
\centering
\begin{tikzpicture}
\begin{axis}[
    width=0.86\textwidth,
    height=0.40\textwidth,
    xlabel={轴向位置 $z$ (\si{\centi\meter})},
    ylabel={计数 $\bar{N}$ (10 s)},
    xmin=0, xmax=38,
    ymin=580, ymax=810,
    grid=both,
    mark size=1.5pt
]
\addplot[blue, thick, mark=*] coordinates {(0,658.0) (1,754.0) (2,763.5) (3,774.5) (4,782.0) (5,786.5) (6,789.5) (7,792.0) (8,790.0) (9,793.5) (10,793.0) (11,794.5) (12,795.0) (13,794.0) (14,794.5) (15,794.0) (16,793.5) (17,793.0) (18,794.5) (19,794.0) (20,792.5) (21,794.5) (22,794.5) (23,791.5) (24,793.0) (25,791.5) (26,791.0) (27,787.0) (28,785.5) (29,784.0) (30,783.5) (31,780.5) (32,776.5) (33,774.0) (34,766.5) (35,755.5) (36,735.0) (37,699.5) (38,608.0)};
\end{axis}
\end{tikzpicture}
\caption{电离室轴向响应曲线（两组均值）\label{fig:axial}}
\end{figure}

\subsection{重复性与误差讨论}
对全部重复测点比较，除低电压小计数点外，绝大多数点两组相对差异在 1\% 以内；说明重复性总体较好。主要误差来源包括：
\begin{enumerate}
    \item 统计涨落（尤其低计数区）；
    \item 高压稳定度与电子学零点漂移；
    \item 准直器与源位置重复定位误差（影响轴向曲线边缘）；
    \item 个别记录异常值（本次已单独标注）。
\end{enumerate}

\section{实验结论}
\begin{enumerate}
    \item 在不同辐照强度和不同压力下，均获得了典型的“上升段 + 饱和平台”V-I 曲线。
    \item 辐照强度减弱（0块板 $\to$ 3块板）时，平台计数显著下降；压力升高（1 MPa $\to$ 2.5 MPa）时，平台计数显著增加。
    \item 结合全部曲线，\SI{300}{\volt} 处处于稳定平台区，可作为本组数据的合适工作电压。
    \item 轴向响应在 \SIrange{4}{31}{cm} 区间内保持较好一致性，非均匀度约 1.83\%，两端响应明显下降。
\end{enumerate}

\section{思考题}
\begin{enumerate}
    \item \textbf{能否使一个电离室既可用来探测 $\beta$ 粒子，又可用于测量 $\gamma$ 射线？}\\
    可以，但需要在结构与工况上折中设计。测 $\beta$ 时通常要求薄窗（减小自吸收）和较浅灵敏区；测 $\gamma$ 时常需要较大体积、适当壁厚（通过次级电子引发电离）以及更高气压以提高灵敏度。实验上可通过更换窗/壁结构或使用不同探测头来分别优化两类测量。

    \item \textbf{若单个辐射粒子能量为 $E$，每对离子产生能量为 $\epsilon$，单位时间有 $n$ 个粒子穿过灵敏体积且能量全部耗尽，电离电流表达式是什么？}\\
    单位时间产生离子对数为 $nE/\epsilon$，故电离电流为
    \begin{equation}
        I = n\frac{E}{\epsilon}\,\mathrm{e}.
    \end{equation}

    \item \textbf{相同活度的 \SI{0.5}{MeV} $\beta$ 源与 \SI{5}{MeV} $\alpha$ 源，在同一空气电离室中测得电流如何比较？设 $\epsilon=\SI{33.85}{eV}$。}\\
    设源活度（单位时间进入灵敏体积的粒子数）为 $A$，则
    \begin{align}
        I_{\beta} &= A\frac{0.5\times10^6}{33.85}\,\mathrm{e}
        \approx 2.37\times10^{-15}A\quad (\si{\ampere}),\\
        I_{\alpha} &= A\frac{5\times10^6}{33.85}\,\mathrm{e}
        \approx 2.37\times10^{-14}A\quad (\si{\ampere}).
    \end{align}
    因此
    \begin{equation}
        I_{\alpha} \approx 10 I_{\beta}.
    \end{equation}
    例如取 $A=\SI{3.7e4}{\becquerel}$，则 $I_{\beta}\approx\SI{8.76e-11}{\ampere}$，$I_{\alpha}\approx\SI{8.76e-10}{\ampere}$。
\end{enumerate}

\end{document}
